\documentclass[10pt,a4paper]{amsart}
\usepackage{amsmath,amssymb,amsthm,graphicx}
% Retained minimal preamble for the Stacks Project chapter sites-cohomology.tex (Cohomology on Sites).
% Verbatim excerpt of the chapter preamble: Xy-pic setup (native preamble lines 14-18) and the native
% theorem declarations with the mathematical macros (native preamble lines 154-217). Copied mechanically;
% nothing added, removed or redefined. The wrapper class and the amsmath/amssymb/amsthm/graphicx base
% packages are supplied by the assembler, so the class/externaldocument part of the native preamble is omitted.
% For commutative diagrams we use Xy-pic
\usepackage[all]{xy}

% We use 2cell for 2-commutative diagrams.
\xyoption{2cell}
% Theorem environments.
%
\theoremstyle{plain}
\newtheorem{theorem}[subsection]{Theorem}
\newtheorem{proposition}[subsection]{Proposition}
\newtheorem{lemma}[subsection]{Lemma}

\theoremstyle{definition}
\newtheorem{definition}[subsection]{Definition}
\newtheorem{example}[subsection]{Example}
\newtheorem{exercise}[subsection]{Exercise}
\newtheorem{situation}[subsection]{Situation}

\theoremstyle{remark}
\newtheorem{remark}[subsection]{Remark}
\newtheorem{remarks}[subsection]{Remarks}

\numberwithin{equation}{subsection}

% Macros
%
\def\lim{\mathop{\mathrm{lim}}\nolimits}
\def\colim{\mathop{\mathrm{colim}}\nolimits}
\def\Spec{\mathop{\mathrm{Spec}}}
\def\Hom{\mathop{\mathrm{Hom}}\nolimits}
\def\Ext{\mathop{\mathrm{Ext}}\nolimits}
\def\SheafHom{\mathop{\mathcal{H}\!\mathit{om}}\nolimits}
\def\SheafExt{\mathop{\mathcal{E}\!\mathit{xt}}\nolimits}
\def\Sch{\mathit{Sch}}
\def\Mor{\mathop{\mathrm{Mor}}\nolimits}
\def\Ob{\mathop{\mathrm{Ob}}\nolimits}
\def\Sh{\mathop{\mathit{Sh}}\nolimits}
\def\NL{\mathop{N\!L}\nolimits}
\def\CH{\mathop{\mathrm{CH}}\nolimits}
\def\proetale{{pro\text{-}\acute{e}tale}}
\def\etale{{\acute{e}tale}}
\def\QCoh{\mathit{QCoh}}
\def\Ker{\mathop{\mathrm{Ker}}}
\def\Im{\mathop{\mathrm{Im}}}
\def\Coker{\mathop{\mathrm{Coker}}}
\def\Coim{\mathop{\mathrm{Coim}}}

% Boxtimes
%
\DeclareMathSymbol{\boxtimes}{\mathbin}{AMSa}{"02}

%
% Macros for moduli stacks/spaces
%
\def\QCohstack{\mathcal{QC}\!\mathit{oh}}
\def\Cohstack{\mathcal{C}\!\mathit{oh}}
\def\Spacesstack{\mathcal{S}\!\mathit{paces}}
\def\Quotfunctor{\mathrm{Quot}}
\def\Hilbfunctor{\mathrm{Hilb}}
\def\Curvesstack{\mathcal{C}\!\mathit{urves}}
\def\Polarizedstack{\mathcal{P}\!\mathit{olarized}}
\def\Complexesstack{\mathcal{C}\!\mathit{omplexes}}
% \Pic is the operator that assigns to X its picard group, usage \Pic(X)
% \Picardstack_{X/B} denotes the Picard stack of X over B
% \Picardfunctor_{X/B} denotes the Picard functor of X over B
\def\Pic{\mathop{\mathrm{Pic}}\nolimits}
\def\Picardstack{\mathcal{P}\!\mathit{ic}}
\def\Picardfunctor{\mathrm{Pic}}
\def\Deformationcategory{\mathcal{D}\!\mathit{ef}}

% --- appended 2026-10-12 (Chief): the project own definitions of the reference/slogan/history
% display environments, copied mechanically from the shared Stacks preamble lines 8-12
% (\usepackage{verbatim} plus the three \newenvironment lines, which the chapter preamble keeps
% verbatim from the project preamble). Nothing added, removed or redefined: these environments
% carry no mathematical content and the project itself typesets them as comments. Without these
% three lines a counted statement whose own text uses them fails with Environment slogan undefined.
\usepackage{verbatim}
\newenvironment{reference}{\comment}{\endcomment}
\newenvironment{slogan}{\comment}{\endcomment}
\newenvironment{history}{\comment}{\endcomment}
\usepackage{cleveref}
\crefname{originallemma}{original lemma}{original lemmas}
\Crefname{originallemma}{Original lemma}{Original lemmas}
\makeatletter
\newlabel{homology-lemma-characterize-weak-serre-subcategory}{{0754}{source}{}{sourcepointer-1}{}}
\newlabel{homology-lemma-characterize-weak-serre-subcategory@cref}{{[originallemma][0][]0754}{[1][0][]source}{}{}{}}
\makeatother
\crefname{originalreference}{original reference}{original references}
\Crefname{originalreference}{Original reference}{Original references}
\makeatletter
\newlabel{item-A-and-P}{{0EZ7}{source}{}{sourcepointer-2}{}}
\newlabel{item-A-and-P@cref}{{[originalreference][0][]0EZ7}{[1][0][]source}{}{}{}}
\makeatother
\crefname{originalreference}{original reference}{original references}
\Crefname{originalreference}{Original reference}{Original references}
\makeatletter
\newlabel{item-A-sheaf}{{0EZ6}{source}{}{sourcepointer-3}{}}
\newlabel{item-A-sheaf@cref}{{[originalreference][0][]0EZ6}{[1][0][]source}{}{}{}}
\makeatother
\crefname{originalreference}{original reference}{original references}
\Crefname{originalreference}{Original reference}{Original references}
\makeatletter
\newlabel{item-base-change-P}{{0EZ4}{source}{}{sourcepointer-4}{}}
\newlabel{item-base-change-P@cref}{{[originalreference][0][]0EZ4}{[1][0][]source}{}{}{}}
\makeatother
\crefname{originalreference}{original reference}{original references}
\Crefname{originalreference}{Original reference}{Original references}
\makeatletter
\newlabel{item-refine-tau-by-P}{{0EZ8}{source}{}{sourcepointer-5}{}}
\newlabel{item-refine-tau-by-P@cref}{{[originalreference][0][]0EZ8}{[1][0][]source}{}{}{}}
\makeatother
\crefname{originalreference}{original reference}{original references}
\Crefname{originalreference}{Original reference}{Original references}
\makeatletter
\newlabel{item-restriction-A}{{0EZ5}{source}{}{sourcepointer-6}{}}
\newlabel{item-restriction-A@cref}{{[originalreference][0][]0EZ5}{[1][0][]source}{}{}{}}
\makeatother
\crefname{originalreference}{original reference}{original references}
\Crefname{originalreference}{Original reference}{Original references}
\makeatletter
\newlabel{lemma-LC-basic}{{09WZ}{source}{}{sourcepointer-7}{}}
\newlabel{lemma-LC-basic@cref}{{[originalreference][0][]09WZ}{[1][0][]source}{}{}{}}
\makeatother
\crefname{originalreference}{original reference}{original references}
\Crefname{originalreference}{Original reference}{Original references}
\makeatletter
\newlabel{lemma-collect-true-things-Zar}{{0DCU}{source}{}{sourcepointer-8}{}}
\newlabel{lemma-collect-true-things-Zar@cref}{{[originalreference][0][]0DCU}{[1][0][]source}{}{}{}}
\makeatother
\crefname{originalreference}{original reference}{original references}
\Crefname{originalreference}{Original reference}{Original references}
\makeatletter
\newlabel{lemma-describe-pullback-pi}{{09X3}{source}{}{sourcepointer-9}{}}
\newlabel{lemma-describe-pullback-pi@cref}{{[originalreference][0][]09X3}{[1][0][]source}{}{}{}}
\makeatother
\crefname{originalreference}{original reference}{original references}
\Crefname{originalreference}{Original reference}{Original references}
\makeatletter
\newlabel{lemma-proper-surjective-is-qc-covering}{{09X5}{source}{}{sourcepointer-10}{}}
\newlabel{lemma-proper-surjective-is-qc-covering@cref}{{[originalreference][0][]09X5}{[1][0][]source}{}{}{}}
\makeatother
\crefname{originalreference}{original reference}{original references}
\Crefname{originalreference}{Original reference}{Original references}
\makeatletter
\newlabel{situation-compare}{{0EZ3}{source}{}{sourcepointer-11}{}}
\newlabel{situation-compare@cref}{{[originalreference][0][]0EZ3}{[1][0][]source}{}{}{}}
\makeatother
\crefname{originallemma}{original lemma}{original lemmas}
\Crefname{originallemma}{Original lemma}{Original lemmas}
\makeatletter
\newlabel{topology-lemma-base-change-separated}{{0CY4}{source}{}{sourcepointer-12}{}}
\newlabel{topology-lemma-base-change-separated@cref}{{[originallemma][0][]0CY4}{[1][0][]source}{}{}{}}
\makeatother
\crefname{originallemma}{original lemma}{original lemmas}
\Crefname{originallemma}{Original lemma}{Original lemmas}
\makeatletter
\newlabel{topology-lemma-closed-map}{{08YD}{source}{}{sourcepointer-13}{}}
\newlabel{topology-lemma-closed-map@cref}{{[originallemma][0][]08YD}{[1][0][]source}{}{}{}}
\makeatother
\crefname{originaltheorem}{original theorem}{original theorems}
\Crefname{originaltheorem}{Original theorem}{Original theorems}
\makeatletter
\newlabel{topology-theorem-characterize-proper}{{005R}{source}{}{sourcepointer-14}{}}
\newlabel{topology-theorem-characterize-proper@cref}{{[originaltheorem][0][]005R}{[1][0][]source}{}{}{}}
\makeatother
\title{The five standing properties hold for the comparison of the locally countably generated site with its Zariski counterpart}
\author{The Stacks Project Authors}
\date{stacks-project:sites-cohomology.tex:local-snapshot-20261011; GNU FDL 1.2 or later}
\begin{document}
\maketitle
\noindent\textit{Editorial note.} All delivered mathematical text is the author's own, copied from the pinned Stacks Project chapter file sites-cohomology.tex; no statement or proof step was rewritten, reordered or removed.\par\medskip
\noindent\textit{Editorial note.} The counted proof is the author's complete argument: the weak Serre verification through the four conditions, the four cited properties, and the explicit construction refining a qc covering by a surjective proper map.\par\medskip
\noindent\textit{Editorial note.} The label lemma-compare-qc-zar names the counted statement, and the proof environment immediately following it in the source is delivered verbatim as its proof.\par\medskip
\noindent\textit{Editorial note.} The counted statement is carried by the author's own lemma environment in the chapter source.\par\medskip
\noindent\textit{Editorial note.} Pre-open gate: python operations/agent-authoring/tools/precard\_gate.py --source sites-cohomology.tex --label lemma-compare-qc-zar --start 7985 --end 8046 --local-id 3556 returned RESULT=OK (L3556 free, occupancy FREE, public interval clear). Candidate taken from lane-s4/halfB-priority.json.\par\medskip
\noindent\textit{Editorial note.} The statement refers to the five labelled items of the standing comparison situation; these are kept as external pointers with their canonical Stacks tags.\par\medskip
\noindent\textit{Editorial note.} The statement and proof are copied from the pinned chapter file.\par\medskip
\noindent\textit{Allowed standard prerequisites (editorial).}\begin{itemize}
\item Characterization of weak Serre subcategories. Homology, the lemma listing conditions (1)-(4) characterizing a weak Serre subcategory of an abelian category. Supplies the four conditions that are checked to show that the full subcategory of sheaves of the form pi\_X\textasciicircum{}\{-1\}F is a weak Serre subcategory.
\item The pullback of the comparison is exact and fully faithful. Same chapter, the lemma collecting the true statements about the Zariski comparison, in particular the exactness and full faithfulness of pi\_X\textasciicircum{}\{-1\} and the exactness of pi\_\{X,*\} on short exact sequences of sheaves in the image. Supplies conditions (1)-(3) of the weak Serre criterion and the pushforward exactness used to check the fourth condition.
\item Basic properties of the locally countably generated site and properness of base changes. Same chapter, the lemma recording the basic properties of the site, together with the topological characterization of proper maps and the base change of separated maps. Supplies property (item-base-change-P), because base change of a proper map is proper.
\item The sheaf property for pullbacks along the projection and the compatibility of the class with proper maps. Same chapter, the lemma describing the pullback along the projection of the site, and part (7)(b) of the lemma collecting true statements about the Zariski comparison. Supply property (item-A-sheaf) and property (item-A-and-P) respectively.
\item A surjective proper map is a qc covering. Same chapter, the lemma that a surjective proper morphism of topological spaces is a qc covering. Supplies the last step of the construction refining a qc covering by a proper map.
\end{itemize}\medskip
\ifdefined\hypertarget\hypertarget{sourcepointer-1}{}\fi
\noindent\textit{Source reference.} Original lemma 0754: sources/stacks-chow/tags.txt, source lines 9118--9118. Prerequisite pointer invoked at proof line(s) 8006; it names a standard fact of a sibling chapter that the counted argument uses, and every step of the counted proof built on it is retained verbatim in the delivered body.\par\medskip
\ifdefined\hypertarget\hypertarget{sourcepointer-2}{}\fi
\noindent\textit{Source reference.} Original reference 0EZ7: sources/stacks-sites-cohomology/pinned-original-source.tex, source lines 6862--6862. Prerequisite pointer invoked at proof line(s) 7998; it names a standard fact of the same chapter that the counted argument uses, and every step of the counted proof built on it is retained verbatim in the delivered body.\par\medskip
\ifdefined\hypertarget\hypertarget{sourcepointer-3}{}\fi
\noindent\textit{Source reference.} Original reference 0EZ6: sources/stacks-sites-cohomology/pinned-original-source.tex, source lines 6857--6857. Prerequisite pointer invoked at proof line(s) 7997; it names a standard fact of the same chapter that the counted argument uses, and every step of the counted proof built on it is retained verbatim in the delivered body.\par\medskip
\ifdefined\hypertarget\hypertarget{sourcepointer-4}{}\fi
\noindent\textit{Source reference.} Original reference 0EZ4: sources/stacks-sites-cohomology/pinned-original-source.tex, source lines 6849--6849. Prerequisite pointer invoked at proof line(s) 7995; it names a standard fact of the same chapter that the counted argument uses, and every step of the counted proof built on it is retained verbatim in the delivered body.\par\medskip
\ifdefined\hypertarget\hypertarget{sourcepointer-5}{}\fi
\noindent\textit{Source reference.} Original reference 0EZ8: sources/stacks-sites-cohomology/pinned-original-source.tex, source lines 6868--6868. Prerequisite pointer invoked at proof line(s) 7999; it names a standard fact of the same chapter that the counted argument uses, and every step of the counted proof built on it is retained verbatim in the delivered body.\par\medskip
\ifdefined\hypertarget\hypertarget{sourcepointer-6}{}\fi
\noindent\textit{Source reference.} Original reference 0EZ5: sources/stacks-sites-cohomology/pinned-original-source.tex, source lines 6853--6853. Prerequisite pointer invoked at proof line(s) 7996; it names a standard fact of the same chapter that the counted argument uses, and every step of the counted proof built on it is retained verbatim in the delivered body.\par\medskip
\ifdefined\hypertarget\hypertarget{sourcepointer-7}{}\fi
\noindent\textit{Source reference.} Original reference 09WZ: sources/stacks-sites-cohomology/pinned-original-source.tex, source lines 7474--7474. Prerequisite pointer invoked at proof line(s) 8017; it names a standard fact of the same chapter that the counted argument uses, and every step of the counted proof built on it is retained verbatim in the delivered body.\par\medskip
\ifdefined\hypertarget\hypertarget{sourcepointer-8}{}\fi
\noindent\textit{Source reference.} Original reference 0DCU: sources/stacks-sites-cohomology/pinned-original-source.tex, source lines 7770--7770. Prerequisite pointer invoked at proof line(s) 8008; it names a standard fact of the same chapter that the counted argument uses, and every step of the counted proof built on it is retained verbatim in the delivered body.\par\medskip
\ifdefined\hypertarget\hypertarget{sourcepointer-9}{}\fi
\noindent\textit{Source reference.} Original reference 09X3: sources/stacks-sites-cohomology/pinned-original-source.tex, source lines 7675--7675. Prerequisite pointer invoked at proof line(s) 8027; it names a standard fact of the same chapter that the counted argument uses, and every step of the counted proof built on it is retained verbatim in the delivered body.\par\medskip
\ifdefined\hypertarget\hypertarget{sourcepointer-10}{}\fi
\noindent\textit{Source reference.} Original reference 09X5: sources/stacks-sites-cohomology/pinned-original-source.tex, source lines 7585--7585. Prerequisite pointer invoked at proof line(s) 8043; it names a standard fact of the same chapter that the counted argument uses, and every step of the counted proof built on it is retained verbatim in the delivered body.\par\medskip
\ifdefined\hypertarget\hypertarget{sourcepointer-11}{}\fi
\noindent\textit{Source reference.} Original reference 0EZ3: sources/stacks-sites-cohomology/pinned-original-source.tex, source lines 6841--6841. Prerequisite pointer invoked at proof line(s) 8000; it names a standard fact of the same chapter that the counted argument uses, and every step of the counted proof built on it is retained verbatim in the delivered body.\par\medskip
\ifdefined\hypertarget\hypertarget{sourcepointer-12}{}\fi
\noindent\textit{Source reference.} Original lemma 0CY4: sources/stacks-chow/tags.txt, source lines 16229--16229. Prerequisite pointer invoked at proof line(s) 8020; it names a standard fact of a sibling chapter that the counted argument uses, and every step of the counted proof built on it is retained verbatim in the delivered body.\par\medskip
\ifdefined\hypertarget\hypertarget{sourcepointer-13}{}\fi
\noindent\textit{Source reference.} Original lemma 08YD: sources/stacks-chow/tags.txt, source lines 11336--11336. Prerequisite pointer invoked at proof line(s) 8040; it names a standard fact of a sibling chapter that the counted argument uses, and every step of the counted proof built on it is retained verbatim in the delivered body.\par\medskip
\ifdefined\hypertarget\hypertarget{sourcepointer-14}{}\fi
\noindent\textit{Source reference.} Original theorem 005R: sources/stacks-chow/tags.txt, source lines 219--219. Prerequisite pointer invoked at proof line(s) 8019; it names a standard fact of a sibling chapter that the counted argument uses, and every step of the counted proof built on it is retained verbatim in the delivered body.\par\medskip
\begin{lemma}
\label{lemma-compare-qc-zar}
Consider the comparison morphism
$\epsilon : \textit{LC}_{qc} \to \textit{LC}_{Zar}$.
Let $\mathcal{P}$ denote the class of proper maps of topological spaces.
For $X$ in $\textit{LC}_{Zar}$ denote
$\mathcal{A}'_X \subset \textit{Ab}(\textit{LC}_{Zar}/X)$
the full subcategory consisting of sheaves of the form
$\pi_X^{-1}\mathcal{F}$ with $\mathcal{F}$ in $\textit{Ab}(X)$.
Then
(\ref{item-base-change-P}),
(\ref{item-restriction-A}),
(\ref{item-A-sheaf}),
(\ref{item-A-and-P}), and
(\ref{item-refine-tau-by-P})
of Situation \ref{situation-compare} hold.
\end{lemma}

\begin{proof}
We first show that $\mathcal{A}'_X \subset \textit{Ab}(\textit{LC}_{Zar}/X)$
is a weak Serre subcategory by checking conditions (1), (2), (3), and (4)
of Homology, Lemma \ref{homology-lemma-characterize-weak-serre-subcategory}.
Parts (1), (2), (3) are immediate as $\pi_X^{-1}$ is exact and
fully faithful by Lemma \ref{lemma-collect-true-things-Zar} part (3). If
$0 \to \pi_X^{-1}\mathcal{F} \to \mathcal{G} \to \pi_X^{-1}\mathcal{F}' \to 0$
is a short exact sequence in $\textit{Ab}(\textit{LC}_{Zar}/X)$
then $0 \to \mathcal{F} \to \pi_{X, *}\mathcal{G} \to \mathcal{F}' \to 0$
is exact by Lemma \ref{lemma-collect-true-things-Zar} part (2).
Hence $\mathcal{G} = \pi_X^{-1}\pi_{X, *}\mathcal{G}$ is in
$\mathcal{A}'_X$ which checks the final condition.

\medskip\noindent
Property (\ref{item-base-change-P}) holds by Lemma \ref{lemma-LC-basic}
and the fact that the base change of a proper map is a proper map (see
Topology, Theorem \ref{topology-theorem-characterize-proper} and
Lemma \ref{topology-lemma-base-change-separated}).

\medskip\noindent
Property (\ref{item-restriction-A}) follows from the commutative
diagram (5) in Lemma \ref{lemma-collect-true-things-Zar}.

\medskip\noindent
Property (\ref{item-A-sheaf}) is Lemma \ref{lemma-describe-pullback-pi}.

\medskip\noindent
Property (\ref{item-A-and-P}) is Lemma \ref{lemma-collect-true-things-Zar}
part (7)(b).

\medskip\noindent
Proof of (\ref{item-refine-tau-by-P}). Suppose given a qc covering
$\{U_i \to U\}$. For $u \in U$ pick $i_1, \ldots, i_m \in I$ and
quasi-compact subsets $E_j \subset U_{i_j}$ such that
$\bigcup f_{i_j}(E_j)$ is a neighbourhood of $u$.
Observe that $Y = \coprod_{j = 1, \ldots, m} E_j \to U$
is proper as a continuous map between Hausdorff quasi-compact spaces
(Topology, Lemma \ref{topology-lemma-closed-map}).
Choose an open neighbourhood $u \in V$ contained in $\bigcup f_{i_j}(E_j)$.
Then $Y \times_U V \to V$ is a surjective proper morphism and
hence a $qc$ covering by Lemma \ref{lemma-proper-surjective-is-qc-covering}.
Since we can do this for every $u \in U$ we see that
(\ref{item-refine-tau-by-P}) holds.
\end{proof}

\end{document}
